% ============================================
% Shared preamble
% ============================================
% Usage: % ============================================
% Shared preamble
% ============================================
% Usage: % ============================================
% Shared preamble
% ============================================
% Usage: % ============================================
% Shared preamble
% ============================================
% Usage: \input{preamble} after \documentclass

% ENCODING & FONTS
\usepackage[utf8]{inputenc}
\usepackage[T1]{fontenc}
\usepackage{lmodern}

% PAGE LAYOUT (Golden Ratio)
\usepackage[margin=1.618cm, top=2.618cm, bottom=2.618cm]{geometry}

% ESSENTIAL PACKAGES
\usepackage{float}
\usepackage{caption}
\usepackage{subcaption}
\usepackage{setspace}
\usepackage{fancyhdr}
\usepackage{xcolor}
\usepackage{hyperref}
\usepackage{csquotes}
\usepackage{amsmath}
\usepackage{amssymb}
\usepackage{amsthm}
\usepackage{mathtools}
\usepackage{booktabs}
\usepackage{longtable}
\usepackage{array}
\usepackage{multirow}
\usepackage{tikz}
\usepackage{graphicx}
\usepackage{listings}
\usepackage{enumitem}
\usepackage{etoolbox}
% Render generated tables a notch smaller so wide pipe-tables breathe.
\AtBeginEnvironment{longtable}{\footnotesize}
\AtBeginEnvironment{tabular}{\small}
% Unicode fallbacks (pandoc artifacts that survive markdown→LaTeX)
\DeclareUnicodeCharacter{00B0}{\ensuremath{^\circ}}
\DeclareUnicodeCharacter{221A}{\ensuremath{\sqrt{\,}}}
\DeclareUnicodeCharacter{2713}{\ensuremath{\checkmark}}
\DeclareUnicodeCharacter{2717}{\ensuremath{\times}}
\DeclareUnicodeCharacter{2070}{\ensuremath{^{0}}}
\DeclareUnicodeCharacter{00B9}{\ensuremath{^{1}}}
\DeclareUnicodeCharacter{00B2}{\ensuremath{^{2}}}
\DeclareUnicodeCharacter{00B3}{\ensuremath{^{3}}}
\DeclareUnicodeCharacter{2074}{\ensuremath{^{4}}}
\DeclareUnicodeCharacter{2075}{\ensuremath{^{5}}}
\DeclareUnicodeCharacter{2076}{\ensuremath{^{6}}}
\DeclareUnicodeCharacter{2077}{\ensuremath{^{7}}}
\DeclareUnicodeCharacter{2078}{\ensuremath{^{8}}}
\DeclareUnicodeCharacter{2079}{\ensuremath{^{9}}}
\DeclareUnicodeCharacter{2071}{\ensuremath{^{i}}}
\DeclareUnicodeCharacter{2080}{\ensuremath{_{0}}}
\DeclareUnicodeCharacter{2081}{\ensuremath{_{1}}}
\DeclareUnicodeCharacter{2082}{\ensuremath{_{2}}}
\DeclareUnicodeCharacter{2083}{\ensuremath{_{3}}}
\DeclareUnicodeCharacter{2084}{\ensuremath{_{4}}}
\DeclareUnicodeCharacter{2085}{\ensuremath{_{5}}}
\DeclareUnicodeCharacter{2086}{\ensuremath{_{6}}}
\DeclareUnicodeCharacter{2087}{\ensuremath{_{7}}}
\DeclareUnicodeCharacter{2088}{\ensuremath{_{8}}}
\DeclareUnicodeCharacter{2089}{\ensuremath{_{9}}}
\DeclareUnicodeCharacter{208A}{\ensuremath{_{+}}}
\DeclareUnicodeCharacter{208B}{\ensuremath{_{-}}}
\DeclareUnicodeCharacter{2192}{\ensuremath{\to}}
\DeclareUnicodeCharacter{2502}{|}
\DeclareUnicodeCharacter{25BC}{\ensuremath{\blacktriangledown}}
\DeclareUnicodeCharacter{2500}{\ensuremath{-}}
\DeclareUnicodeCharacter{2514}{\ensuremath{\llcorner}}
\DeclareUnicodeCharacter{251C}{\ensuremath{\vdash}}
\DeclareUnicodeCharacter{2026}{\ldots}
\DeclareUnicodeCharacter{2208}{\ensuremath{\in}}
\DeclareUnicodeCharacter{2229}{\ensuremath{\cap}}
\DeclareUnicodeCharacter{222A}{\ensuremath{\cup}}
\DeclareUnicodeCharacter{2295}{\ensuremath{\oplus}}
\DeclareUnicodeCharacter{2297}{\ensuremath{\otimes}}
\DeclareUnicodeCharacter{2245}{\ensuremath{\cong}}
\DeclareUnicodeCharacter{2248}{\ensuremath{\approx}}
\DeclareUnicodeCharacter{2260}{\ensuremath{\neq}}
\DeclareUnicodeCharacter{2261}{\ensuremath{\equiv}}
\DeclareUnicodeCharacter{2264}{\ensuremath{\leq}}
\DeclareUnicodeCharacter{2265}{\ensuremath{\geq}}
\DeclareUnicodeCharacter{00D7}{\ensuremath{\times}}
\DeclareUnicodeCharacter{00B1}{\ensuremath{\pm}}
\DeclareUnicodeCharacter{2225}{\ensuremath{\|}}
% Greek letters used in body text (pandoc passes them through unchanged)
\DeclareUnicodeCharacter{03B1}{\ensuremath{\alpha}}
\DeclareUnicodeCharacter{03B2}{\ensuremath{\beta}}
\DeclareUnicodeCharacter{03B3}{\ensuremath{\gamma}}
\DeclareUnicodeCharacter{03B4}{\ensuremath{\delta}}
\DeclareUnicodeCharacter{03B5}{\ensuremath{\varepsilon}}
\DeclareUnicodeCharacter{03B6}{\ensuremath{\zeta}}
\DeclareUnicodeCharacter{03B7}{\ensuremath{\eta}}
\DeclareUnicodeCharacter{03B8}{\ensuremath{\theta}}
\DeclareUnicodeCharacter{03B9}{\ensuremath{\iota}}
\DeclareUnicodeCharacter{03BA}{\ensuremath{\kappa}}
\DeclareUnicodeCharacter{03BB}{\ensuremath{\lambda}}
\DeclareUnicodeCharacter{03BC}{\ensuremath{\mu}}
\DeclareUnicodeCharacter{03BD}{\ensuremath{\nu}}
\DeclareUnicodeCharacter{03BE}{\ensuremath{\xi}}
\DeclareUnicodeCharacter{03C0}{\ensuremath{\pi}}
\DeclareUnicodeCharacter{03C1}{\ensuremath{\rho}}
\DeclareUnicodeCharacter{03C3}{\ensuremath{\sigma}}
\DeclareUnicodeCharacter{03C4}{\ensuremath{\tau}}
\DeclareUnicodeCharacter{03C5}{\ensuremath{\upsilon}}
\DeclareUnicodeCharacter{03C6}{\ensuremath{\varphi}}
\DeclareUnicodeCharacter{03C7}{\ensuremath{\chi}}
\DeclareUnicodeCharacter{03C8}{\ensuremath{\psi}}
\DeclareUnicodeCharacter{03C9}{\ensuremath{\omega}}
\DeclareUnicodeCharacter{0394}{\ensuremath{\Delta}}
\DeclareUnicodeCharacter{0398}{\ensuremath{\Theta}}
\DeclareUnicodeCharacter{039B}{\ensuremath{\Lambda}}
\DeclareUnicodeCharacter{03A0}{\ensuremath{\Pi}}
\DeclareUnicodeCharacter{03A3}{\ensuremath{\Sigma}}
\DeclareUnicodeCharacter{03A6}{\ensuremath{\Phi}}
\DeclareUnicodeCharacter{03A8}{\ensuremath{\Psi}}
\DeclareUnicodeCharacter{03A9}{\ensuremath{\Omega}}
\DeclareUnicodeCharacter{2200}{\ensuremath{\forall}}
\DeclareUnicodeCharacter{2203}{\ensuremath{\exists}}
\DeclareUnicodeCharacter{2207}{\ensuremath{\nabla}}
\DeclareUnicodeCharacter{2202}{\ensuremath{\partial}}
\DeclareUnicodeCharacter{2218}{\ensuremath{\circ}}
\DeclareUnicodeCharacter{2032}{\ensuremath{\prime}}
\DeclareUnicodeCharacter{2113}{\ensuremath{\ell}}
\DeclareUnicodeCharacter{210F}{\ensuremath{\hbar}}
\DeclareUnicodeCharacter{27E8}{\ensuremath{\langle}}
\DeclareUnicodeCharacter{27E9}{\ensuremath{\rangle}}
\DeclareUnicodeCharacter{2308}{\ensuremath{\lceil}}
\DeclareUnicodeCharacter{2309}{\ensuremath{\rceil}}
\DeclareUnicodeCharacter{230A}{\ensuremath{\lfloor}}
\DeclareUnicodeCharacter{230B}{\ensuremath{\rfloor}}
\DeclareUnicodeCharacter{222B}{\ensuremath{\int}}
\DeclareUnicodeCharacter{2211}{\ensuremath{\sum}}
\DeclareUnicodeCharacter{220F}{\ensuremath{\prod}}
\DeclareUnicodeCharacter{221E}{\ensuremath{\infty}}
\DeclareUnicodeCharacter{00B5}{\ensuremath{\mu}}
\DeclareUnicodeCharacter{2236}{\ensuremath{\colon}}
\DeclareUnicodeCharacter{2192}{\ensuremath{\to}}
\DeclareUnicodeCharacter{21A6}{\ensuremath{\mapsto}}
\DeclareUnicodeCharacter{21D2}{\ensuremath{\Rightarrow}}
\DeclareUnicodeCharacter{21D4}{\ensuremath{\Leftrightarrow}}
\DeclareUnicodeCharacter{2282}{\ensuremath{\subset}}
\DeclareUnicodeCharacter{2286}{\ensuremath{\subseteq}}
\DeclareUnicodeCharacter{2287}{\ensuremath{\supseteq}}
\DeclareUnicodeCharacter{222B}{\ensuremath{\int}}
\DeclareUnicodeCharacter{2205}{\ensuremath{\emptyset}}
\DeclareUnicodeCharacter{2299}{\ensuremath{\odot}}
\DeclareUnicodeCharacter{2220}{\ensuremath{\angle}}
\DeclareUnicodeCharacter{2135}{\ensuremath{\aleph}}
\DeclareUnicodeCharacter{207B}{\ensuremath{^{-}}}
\DeclareUnicodeCharacter{207A}{\ensuremath{^{+}}}
\DeclareUnicodeCharacter{220E}{\ensuremath{\blacksquare}}
\DeclareUnicodeCharacter{2061}{}
\DeclareUnicodeCharacter{2062}{}
\DeclareUnicodeCharacter{2063}{}
\DeclareUnicodeCharacter{2009}{\,}
\DeclareUnicodeCharacter{200A}{\,}
\DeclareUnicodeCharacter{200B}{}
\DeclareUnicodeCharacter{2212}{\ensuremath{-}}
\DeclareUnicodeCharacter{2194}{\ensuremath{\leftrightarrow}}
\DeclareUnicodeCharacter{2099}{\ensuremath{_{n}}}
\DeclareUnicodeCharacter{2098}{\ensuremath{_{m}}}
\DeclareUnicodeCharacter{2096}{\ensuremath{_{k}}}
\DeclareUnicodeCharacter{2095}{\ensuremath{_{h}}}
\DeclareUnicodeCharacter{2090}{\ensuremath{_{a}}}
\DeclareUnicodeCharacter{2091}{\ensuremath{_{e}}}
\DeclareUnicodeCharacter{2093}{\ensuremath{_{x}}}
\DeclareUnicodeCharacter{1D62}{\ensuremath{_{i}}}
\DeclareUnicodeCharacter{2C7C}{\ensuremath{_{j}}}
\DeclareUnicodeCharacter{2C7B}{\ensuremath{^{e}}}
\DeclareUnicodeCharacter{1D52}{\ensuremath{^{o}}}
\DeclareUnicodeCharacter{1D5}{\ensuremath{^{u}}}
\DeclareUnicodeCharacter{02E1}{\ensuremath{^{l}}}
\DeclareUnicodeCharacter{2295}{\ensuremath{\oplus}}
\DeclareUnicodeCharacter{27FA}{\ensuremath{\Longleftrightarrow}}
\DeclareUnicodeCharacter{27F9}{\ensuremath{\Longrightarrow}}
\DeclareUnicodeCharacter{2299}{\ensuremath{\odot}}
% FoP-specific symbols (blackboard division algebras, set/relation operators)
\DeclareUnicodeCharacter{00F7}{\ensuremath{\div}}
\DeclareUnicodeCharacter{2102}{\ensuremath{\mathbb{C}}}
\DeclareUnicodeCharacter{210D}{\ensuremath{\mathbb{H}}}
\DeclareUnicodeCharacter{211D}{\ensuremath{\mathbb{R}}}
\DeclareUnicodeCharacter{1D546}{\ensuremath{\mathbb{O}}}
\DeclareUnicodeCharacter{2227}{\ensuremath{\wedge}}
\DeclareUnicodeCharacter{2272}{\ensuremath{\lesssim}}
\DeclareUnicodeCharacter{2283}{\ensuremath{\supset}}
\DeclareUnicodeCharacter{228A}{\ensuremath{\subsetneq}}
\DeclareUnicodeCharacter{2016}{\ensuremath{\|}}
\DeclareUnicodeCharacter{2190}{\ensuremath{\leftarrow}}
\DeclareUnicodeCharacter{2228}{\ensuremath{\vee}}
\DeclareUnicodeCharacter{22C6}{\ensuremath{\star}}
\DeclareUnicodeCharacter{1D4D8}{\ensuremath{\mathcal{I}}}
\DeclareUnicodeCharacter{1D524}{\ensuremath{\mathfrak{g}}}
\DeclareUnicodeCharacter{23F3}{\textsf{[open]}}
% Provide \tightlist used by pandoc-generated lists (no-op)
\providecommand{\tightlist}{\setlength{\itemsep}{0pt}\setlength{\parskip}{0pt}}
% pandoc --listings emits \passthrough around inline \lstinline spans
\providecommand{\passthrough}[1]{#1}
% pandoc longtable column specs: provide \real as passthrough and a
% `none` counter so that calc's \real{x} expansion path works.
\newcounter{none}
\providecommand{\real}[1]{#1}
% pandoc >=3.x wraps figures in \pandocbounded (official definition: scale to
% fit text width/height, never upscale). Requires graphicx (loaded above).
\makeatletter
\newsavebox\pandoc@box
\providecommand*\pandocbounded[1]{%
  \sbox\pandoc@box{#1}%
  \ifdim\wd\pandoc@box>\linewidth
    \resizebox{\linewidth}{!}{\usebox\pandoc@box}%
  \else
    \usebox\pandoc@box
  \fi}
\makeatother

% LISTINGS
\lstset{
    basicstyle=\small\ttfamily,
    breaklines=true, frame=single,
    keepspaces=true, showstringspaces=false,
    aboveskip=0.8em, belowskip=0.8em
}

% HEADER/FOOTER
\setlength{\headheight}{14pt}
\pagestyle{fancy}
\fancyhf{}
\fancyhead[R]{\thepage}
\renewcommand{\headrulewidth}{0.2pt}

% HYPERREF
\hypersetup{
    colorlinks=true,
    linkcolor=blue, citecolor=blue, urlcolor=blue,
    pdfauthor={Brieuc de La Fourniere}
}

% SPACING
\setstretch{1.2}
\setlength{\parskip}{0.4em}
\setlength{\parindent}{0pt}

% TITLE FORMATTING
\usepackage{titling}
\pretitle{\LARGE\bfseries}
\posttitle{\vspace{-0.4em}}
\preauthor{}
\postauthor{}
\predate{}
\postdate{}
\setlength{\droptitle}{-2.0em}

% CUSTOM COMMANDS
\newcommand{\E}{\mathrm{E}}
\newcommand{\Gtwo}{\mathrm{G}_2}
\newcommand{\Kseven}{K_7}
\newcommand{\AdS}{\mathrm{AdS}}
\newcommand{\SU}{\mathrm{SU}}
\newcommand{\SO}{\mathrm{SO}}
\newcommand{\U}{\mathrm{U}}
\newcommand{\Sp}{\mathrm{Sp}}
\newcommand{\PSL}{\mathrm{PSL}}
\newcommand{\SM}{\mathrm{SM}}
\newcommand{\Tr}{\mathrm{Tr}}
\newcommand{\rank}{\mathrm{rank}}
\newcommand{\Vol}{\mathrm{Vol}}
\newcommand{\Hol}{\mathrm{Hol}}
\newcommand{\Ric}{\mathrm{Ric}}
\newcommand{\Rm}{\mathrm{Rm}}
\newcommand{\Scal}{\mathrm{Scal}}
\DeclareMathOperator{\diag}{diag}
\DeclareMathOperator{\cond}{cond}
\newcommand{\CP}{\mathrm{CP}}
\newcommand{\CKM}{\mathrm{CKM}}
\newcommand{\PMNS}{\mathrm{PMNS}}
\newcommand{\EW}{\mathrm{EW}}
\newcommand{\Pl}{\mathrm{Pl}}
\newcommand{\DE}{\mathrm{DE}}
\newcommand{\DM}{\mathrm{DM}}
\newcommand{\KK}{\mathrm{KK}}
\newcommand{\NK}{\mathrm{NK}}
\newcommand{\TCS}{\mathrm{TCS}}
\newcommand{\PINN}{\mathrm{PINN}}
\newcommand{\btwo}{b_2}
\newcommand{\bthree}{b_3}
\newcommand{\typeI}{\textsc{I}}
\newcommand{\typeII}{\textsc{II}}
\newcommand{\typeIII}{\textsc{III}}
\newcommand{\typeIV}{\textsc{IV}}
\newcommand{\proven}{\textsc{Proven}}

% PDF bookmark safety
\pdfstringdefDisableCommands{%
  \def\textsubscript#1{#1}%
  \def\textsuperscript#1{#1}%
  \def\CP{CP}%
  \def\Gtwo{G2}%
  \def\E{E}%
  \def\SM{SM}%
}

% THEOREM ENVIRONMENTS
\newtheorem{theorem}{Theorem}[section]
\newtheorem{proposition}[theorem]{Proposition}
\newtheorem{lemma}[theorem]{Lemma}
\newtheorem{corollary}[theorem]{Corollary}
\theoremstyle{definition}
\newtheorem{definition}[theorem]{Definition}
\theoremstyle{remark}
\newtheorem{remark}[theorem]{Remark}
 after \documentclass

% ENCODING & FONTS
\usepackage[utf8]{inputenc}
\usepackage[T1]{fontenc}
\usepackage{lmodern}

% PAGE LAYOUT (Golden Ratio)
\usepackage[margin=1.618cm, top=2.618cm, bottom=2.618cm]{geometry}

% ESSENTIAL PACKAGES
\usepackage{float}
\usepackage{caption}
\usepackage{subcaption}
\usepackage{setspace}
\usepackage{fancyhdr}
\usepackage{xcolor}
\usepackage{hyperref}
\usepackage{csquotes}
\usepackage{amsmath}
\usepackage{amssymb}
\usepackage{amsthm}
\usepackage{mathtools}
\usepackage{booktabs}
\usepackage{longtable}
\usepackage{array}
\usepackage{multirow}
\usepackage{tikz}
\usepackage{graphicx}
\usepackage{listings}
\usepackage{enumitem}
\usepackage{etoolbox}
% Render generated tables a notch smaller so wide pipe-tables breathe.
\AtBeginEnvironment{longtable}{\footnotesize}
\AtBeginEnvironment{tabular}{\small}
% Unicode fallbacks (pandoc artifacts that survive markdown→LaTeX)
\DeclareUnicodeCharacter{00B0}{\ensuremath{^\circ}}
\DeclareUnicodeCharacter{221A}{\ensuremath{\sqrt{\,}}}
\DeclareUnicodeCharacter{2713}{\ensuremath{\checkmark}}
\DeclareUnicodeCharacter{2717}{\ensuremath{\times}}
\DeclareUnicodeCharacter{2070}{\ensuremath{^{0}}}
\DeclareUnicodeCharacter{00B9}{\ensuremath{^{1}}}
\DeclareUnicodeCharacter{00B2}{\ensuremath{^{2}}}
\DeclareUnicodeCharacter{00B3}{\ensuremath{^{3}}}
\DeclareUnicodeCharacter{2074}{\ensuremath{^{4}}}
\DeclareUnicodeCharacter{2075}{\ensuremath{^{5}}}
\DeclareUnicodeCharacter{2076}{\ensuremath{^{6}}}
\DeclareUnicodeCharacter{2077}{\ensuremath{^{7}}}
\DeclareUnicodeCharacter{2078}{\ensuremath{^{8}}}
\DeclareUnicodeCharacter{2079}{\ensuremath{^{9}}}
\DeclareUnicodeCharacter{2071}{\ensuremath{^{i}}}
\DeclareUnicodeCharacter{2080}{\ensuremath{_{0}}}
\DeclareUnicodeCharacter{2081}{\ensuremath{_{1}}}
\DeclareUnicodeCharacter{2082}{\ensuremath{_{2}}}
\DeclareUnicodeCharacter{2083}{\ensuremath{_{3}}}
\DeclareUnicodeCharacter{2084}{\ensuremath{_{4}}}
\DeclareUnicodeCharacter{2085}{\ensuremath{_{5}}}
\DeclareUnicodeCharacter{2086}{\ensuremath{_{6}}}
\DeclareUnicodeCharacter{2087}{\ensuremath{_{7}}}
\DeclareUnicodeCharacter{2088}{\ensuremath{_{8}}}
\DeclareUnicodeCharacter{2089}{\ensuremath{_{9}}}
\DeclareUnicodeCharacter{208A}{\ensuremath{_{+}}}
\DeclareUnicodeCharacter{208B}{\ensuremath{_{-}}}
\DeclareUnicodeCharacter{2192}{\ensuremath{\to}}
\DeclareUnicodeCharacter{2502}{|}
\DeclareUnicodeCharacter{25BC}{\ensuremath{\blacktriangledown}}
\DeclareUnicodeCharacter{2500}{\ensuremath{-}}
\DeclareUnicodeCharacter{2514}{\ensuremath{\llcorner}}
\DeclareUnicodeCharacter{251C}{\ensuremath{\vdash}}
\DeclareUnicodeCharacter{2026}{\ldots}
\DeclareUnicodeCharacter{2208}{\ensuremath{\in}}
\DeclareUnicodeCharacter{2229}{\ensuremath{\cap}}
\DeclareUnicodeCharacter{222A}{\ensuremath{\cup}}
\DeclareUnicodeCharacter{2295}{\ensuremath{\oplus}}
\DeclareUnicodeCharacter{2297}{\ensuremath{\otimes}}
\DeclareUnicodeCharacter{2245}{\ensuremath{\cong}}
\DeclareUnicodeCharacter{2248}{\ensuremath{\approx}}
\DeclareUnicodeCharacter{2260}{\ensuremath{\neq}}
\DeclareUnicodeCharacter{2261}{\ensuremath{\equiv}}
\DeclareUnicodeCharacter{2264}{\ensuremath{\leq}}
\DeclareUnicodeCharacter{2265}{\ensuremath{\geq}}
\DeclareUnicodeCharacter{00D7}{\ensuremath{\times}}
\DeclareUnicodeCharacter{00B1}{\ensuremath{\pm}}
\DeclareUnicodeCharacter{2225}{\ensuremath{\|}}
% Greek letters used in body text (pandoc passes them through unchanged)
\DeclareUnicodeCharacter{03B1}{\ensuremath{\alpha}}
\DeclareUnicodeCharacter{03B2}{\ensuremath{\beta}}
\DeclareUnicodeCharacter{03B3}{\ensuremath{\gamma}}
\DeclareUnicodeCharacter{03B4}{\ensuremath{\delta}}
\DeclareUnicodeCharacter{03B5}{\ensuremath{\varepsilon}}
\DeclareUnicodeCharacter{03B6}{\ensuremath{\zeta}}
\DeclareUnicodeCharacter{03B7}{\ensuremath{\eta}}
\DeclareUnicodeCharacter{03B8}{\ensuremath{\theta}}
\DeclareUnicodeCharacter{03B9}{\ensuremath{\iota}}
\DeclareUnicodeCharacter{03BA}{\ensuremath{\kappa}}
\DeclareUnicodeCharacter{03BB}{\ensuremath{\lambda}}
\DeclareUnicodeCharacter{03BC}{\ensuremath{\mu}}
\DeclareUnicodeCharacter{03BD}{\ensuremath{\nu}}
\DeclareUnicodeCharacter{03BE}{\ensuremath{\xi}}
\DeclareUnicodeCharacter{03C0}{\ensuremath{\pi}}
\DeclareUnicodeCharacter{03C1}{\ensuremath{\rho}}
\DeclareUnicodeCharacter{03C3}{\ensuremath{\sigma}}
\DeclareUnicodeCharacter{03C4}{\ensuremath{\tau}}
\DeclareUnicodeCharacter{03C5}{\ensuremath{\upsilon}}
\DeclareUnicodeCharacter{03C6}{\ensuremath{\varphi}}
\DeclareUnicodeCharacter{03C7}{\ensuremath{\chi}}
\DeclareUnicodeCharacter{03C8}{\ensuremath{\psi}}
\DeclareUnicodeCharacter{03C9}{\ensuremath{\omega}}
\DeclareUnicodeCharacter{0394}{\ensuremath{\Delta}}
\DeclareUnicodeCharacter{0398}{\ensuremath{\Theta}}
\DeclareUnicodeCharacter{039B}{\ensuremath{\Lambda}}
\DeclareUnicodeCharacter{03A0}{\ensuremath{\Pi}}
\DeclareUnicodeCharacter{03A3}{\ensuremath{\Sigma}}
\DeclareUnicodeCharacter{03A6}{\ensuremath{\Phi}}
\DeclareUnicodeCharacter{03A8}{\ensuremath{\Psi}}
\DeclareUnicodeCharacter{03A9}{\ensuremath{\Omega}}
\DeclareUnicodeCharacter{2200}{\ensuremath{\forall}}
\DeclareUnicodeCharacter{2203}{\ensuremath{\exists}}
\DeclareUnicodeCharacter{2207}{\ensuremath{\nabla}}
\DeclareUnicodeCharacter{2202}{\ensuremath{\partial}}
\DeclareUnicodeCharacter{2218}{\ensuremath{\circ}}
\DeclareUnicodeCharacter{2032}{\ensuremath{\prime}}
\DeclareUnicodeCharacter{2113}{\ensuremath{\ell}}
\DeclareUnicodeCharacter{210F}{\ensuremath{\hbar}}
\DeclareUnicodeCharacter{27E8}{\ensuremath{\langle}}
\DeclareUnicodeCharacter{27E9}{\ensuremath{\rangle}}
\DeclareUnicodeCharacter{2308}{\ensuremath{\lceil}}
\DeclareUnicodeCharacter{2309}{\ensuremath{\rceil}}
\DeclareUnicodeCharacter{230A}{\ensuremath{\lfloor}}
\DeclareUnicodeCharacter{230B}{\ensuremath{\rfloor}}
\DeclareUnicodeCharacter{222B}{\ensuremath{\int}}
\DeclareUnicodeCharacter{2211}{\ensuremath{\sum}}
\DeclareUnicodeCharacter{220F}{\ensuremath{\prod}}
\DeclareUnicodeCharacter{221E}{\ensuremath{\infty}}
\DeclareUnicodeCharacter{00B5}{\ensuremath{\mu}}
\DeclareUnicodeCharacter{2236}{\ensuremath{\colon}}
\DeclareUnicodeCharacter{2192}{\ensuremath{\to}}
\DeclareUnicodeCharacter{21A6}{\ensuremath{\mapsto}}
\DeclareUnicodeCharacter{21D2}{\ensuremath{\Rightarrow}}
\DeclareUnicodeCharacter{21D4}{\ensuremath{\Leftrightarrow}}
\DeclareUnicodeCharacter{2282}{\ensuremath{\subset}}
\DeclareUnicodeCharacter{2286}{\ensuremath{\subseteq}}
\DeclareUnicodeCharacter{2287}{\ensuremath{\supseteq}}
\DeclareUnicodeCharacter{222B}{\ensuremath{\int}}
\DeclareUnicodeCharacter{2205}{\ensuremath{\emptyset}}
\DeclareUnicodeCharacter{2299}{\ensuremath{\odot}}
\DeclareUnicodeCharacter{2220}{\ensuremath{\angle}}
\DeclareUnicodeCharacter{2135}{\ensuremath{\aleph}}
\DeclareUnicodeCharacter{207B}{\ensuremath{^{-}}}
\DeclareUnicodeCharacter{207A}{\ensuremath{^{+}}}
\DeclareUnicodeCharacter{220E}{\ensuremath{\blacksquare}}
\DeclareUnicodeCharacter{2061}{}
\DeclareUnicodeCharacter{2062}{}
\DeclareUnicodeCharacter{2063}{}
\DeclareUnicodeCharacter{2009}{\,}
\DeclareUnicodeCharacter{200A}{\,}
\DeclareUnicodeCharacter{200B}{}
\DeclareUnicodeCharacter{2212}{\ensuremath{-}}
\DeclareUnicodeCharacter{2194}{\ensuremath{\leftrightarrow}}
\DeclareUnicodeCharacter{2099}{\ensuremath{_{n}}}
\DeclareUnicodeCharacter{2098}{\ensuremath{_{m}}}
\DeclareUnicodeCharacter{2096}{\ensuremath{_{k}}}
\DeclareUnicodeCharacter{2095}{\ensuremath{_{h}}}
\DeclareUnicodeCharacter{2090}{\ensuremath{_{a}}}
\DeclareUnicodeCharacter{2091}{\ensuremath{_{e}}}
\DeclareUnicodeCharacter{2093}{\ensuremath{_{x}}}
\DeclareUnicodeCharacter{1D62}{\ensuremath{_{i}}}
\DeclareUnicodeCharacter{2C7C}{\ensuremath{_{j}}}
\DeclareUnicodeCharacter{2C7B}{\ensuremath{^{e}}}
\DeclareUnicodeCharacter{1D52}{\ensuremath{^{o}}}
\DeclareUnicodeCharacter{1D5}{\ensuremath{^{u}}}
\DeclareUnicodeCharacter{02E1}{\ensuremath{^{l}}}
\DeclareUnicodeCharacter{2295}{\ensuremath{\oplus}}
\DeclareUnicodeCharacter{27FA}{\ensuremath{\Longleftrightarrow}}
\DeclareUnicodeCharacter{27F9}{\ensuremath{\Longrightarrow}}
\DeclareUnicodeCharacter{2299}{\ensuremath{\odot}}
% FoP-specific symbols (blackboard division algebras, set/relation operators)
\DeclareUnicodeCharacter{00F7}{\ensuremath{\div}}
\DeclareUnicodeCharacter{2102}{\ensuremath{\mathbb{C}}}
\DeclareUnicodeCharacter{210D}{\ensuremath{\mathbb{H}}}
\DeclareUnicodeCharacter{211D}{\ensuremath{\mathbb{R}}}
\DeclareUnicodeCharacter{1D546}{\ensuremath{\mathbb{O}}}
\DeclareUnicodeCharacter{2227}{\ensuremath{\wedge}}
\DeclareUnicodeCharacter{2272}{\ensuremath{\lesssim}}
\DeclareUnicodeCharacter{2283}{\ensuremath{\supset}}
\DeclareUnicodeCharacter{228A}{\ensuremath{\subsetneq}}
\DeclareUnicodeCharacter{2016}{\ensuremath{\|}}
\DeclareUnicodeCharacter{2190}{\ensuremath{\leftarrow}}
\DeclareUnicodeCharacter{2228}{\ensuremath{\vee}}
\DeclareUnicodeCharacter{22C6}{\ensuremath{\star}}
\DeclareUnicodeCharacter{1D4D8}{\ensuremath{\mathcal{I}}}
\DeclareUnicodeCharacter{1D524}{\ensuremath{\mathfrak{g}}}
\DeclareUnicodeCharacter{23F3}{\textsf{[open]}}
% Provide \tightlist used by pandoc-generated lists (no-op)
\providecommand{\tightlist}{\setlength{\itemsep}{0pt}\setlength{\parskip}{0pt}}
% pandoc --listings emits \passthrough around inline \lstinline spans
\providecommand{\passthrough}[1]{#1}
% pandoc longtable column specs: provide \real as passthrough and a
% `none` counter so that calc's \real{x} expansion path works.
\newcounter{none}
\providecommand{\real}[1]{#1}
% pandoc >=3.x wraps figures in \pandocbounded (official definition: scale to
% fit text width/height, never upscale). Requires graphicx (loaded above).
\makeatletter
\newsavebox\pandoc@box
\providecommand*\pandocbounded[1]{%
  \sbox\pandoc@box{#1}%
  \ifdim\wd\pandoc@box>\linewidth
    \resizebox{\linewidth}{!}{\usebox\pandoc@box}%
  \else
    \usebox\pandoc@box
  \fi}
\makeatother

% LISTINGS
\lstset{
    basicstyle=\small\ttfamily,
    breaklines=true, frame=single,
    keepspaces=true, showstringspaces=false,
    aboveskip=0.8em, belowskip=0.8em
}

% HEADER/FOOTER
\setlength{\headheight}{14pt}
\pagestyle{fancy}
\fancyhf{}
\fancyhead[R]{\thepage}
\renewcommand{\headrulewidth}{0.2pt}

% HYPERREF
\hypersetup{
    colorlinks=true,
    linkcolor=blue, citecolor=blue, urlcolor=blue,
    pdfauthor={Brieuc de La Fourniere}
}

% SPACING
\setstretch{1.2}
\setlength{\parskip}{0.4em}
\setlength{\parindent}{0pt}

% TITLE FORMATTING
\usepackage{titling}
\pretitle{\LARGE\bfseries}
\posttitle{\vspace{-0.4em}}
\preauthor{}
\postauthor{}
\predate{}
\postdate{}
\setlength{\droptitle}{-2.0em}

% CUSTOM COMMANDS
\newcommand{\E}{\mathrm{E}}
\newcommand{\Gtwo}{\mathrm{G}_2}
\newcommand{\Kseven}{K_7}
\newcommand{\AdS}{\mathrm{AdS}}
\newcommand{\SU}{\mathrm{SU}}
\newcommand{\SO}{\mathrm{SO}}
\newcommand{\U}{\mathrm{U}}
\newcommand{\Sp}{\mathrm{Sp}}
\newcommand{\PSL}{\mathrm{PSL}}
\newcommand{\SM}{\mathrm{SM}}
\newcommand{\Tr}{\mathrm{Tr}}
\newcommand{\rank}{\mathrm{rank}}
\newcommand{\Vol}{\mathrm{Vol}}
\newcommand{\Hol}{\mathrm{Hol}}
\newcommand{\Ric}{\mathrm{Ric}}
\newcommand{\Rm}{\mathrm{Rm}}
\newcommand{\Scal}{\mathrm{Scal}}
\DeclareMathOperator{\diag}{diag}
\DeclareMathOperator{\cond}{cond}
\newcommand{\CP}{\mathrm{CP}}
\newcommand{\CKM}{\mathrm{CKM}}
\newcommand{\PMNS}{\mathrm{PMNS}}
\newcommand{\EW}{\mathrm{EW}}
\newcommand{\Pl}{\mathrm{Pl}}
\newcommand{\DE}{\mathrm{DE}}
\newcommand{\DM}{\mathrm{DM}}
\newcommand{\KK}{\mathrm{KK}}
\newcommand{\NK}{\mathrm{NK}}
\newcommand{\TCS}{\mathrm{TCS}}
\newcommand{\PINN}{\mathrm{PINN}}
\newcommand{\btwo}{b_2}
\newcommand{\bthree}{b_3}
\newcommand{\typeI}{\textsc{I}}
\newcommand{\typeII}{\textsc{II}}
\newcommand{\typeIII}{\textsc{III}}
\newcommand{\typeIV}{\textsc{IV}}
\newcommand{\proven}{\textsc{Proven}}

% PDF bookmark safety
\pdfstringdefDisableCommands{%
  \def\textsubscript#1{#1}%
  \def\textsuperscript#1{#1}%
  \def\CP{CP}%
  \def\Gtwo{G2}%
  \def\E{E}%
  \def\SM{SM}%
}

% THEOREM ENVIRONMENTS
\newtheorem{theorem}{Theorem}[section]
\newtheorem{proposition}[theorem]{Proposition}
\newtheorem{lemma}[theorem]{Lemma}
\newtheorem{corollary}[theorem]{Corollary}
\theoremstyle{definition}
\newtheorem{definition}[theorem]{Definition}
\theoremstyle{remark}
\newtheorem{remark}[theorem]{Remark}
 after \documentclass

% ENCODING & FONTS
\usepackage[utf8]{inputenc}
\usepackage[T1]{fontenc}
\usepackage{lmodern}

% PAGE LAYOUT (Golden Ratio)
\usepackage[margin=1.618cm, top=2.618cm, bottom=2.618cm]{geometry}

% ESSENTIAL PACKAGES
\usepackage{float}
\usepackage{caption}
\usepackage{subcaption}
\usepackage{setspace}
\usepackage{fancyhdr}
\usepackage{xcolor}
\usepackage{hyperref}
\usepackage{csquotes}
\usepackage{amsmath}
\usepackage{amssymb}
\usepackage{amsthm}
\usepackage{mathtools}
\usepackage{booktabs}
\usepackage{longtable}
\usepackage{array}
\usepackage{multirow}
\usepackage{tikz}
\usepackage{graphicx}
\usepackage{listings}
\usepackage{enumitem}
\usepackage{etoolbox}
% Render generated tables a notch smaller so wide pipe-tables breathe.
\AtBeginEnvironment{longtable}{\footnotesize}
\AtBeginEnvironment{tabular}{\small}
% Unicode fallbacks (pandoc artifacts that survive markdown→LaTeX)
\DeclareUnicodeCharacter{00B0}{\ensuremath{^\circ}}
\DeclareUnicodeCharacter{221A}{\ensuremath{\sqrt{\,}}}
\DeclareUnicodeCharacter{2713}{\ensuremath{\checkmark}}
\DeclareUnicodeCharacter{2717}{\ensuremath{\times}}
\DeclareUnicodeCharacter{2070}{\ensuremath{^{0}}}
\DeclareUnicodeCharacter{00B9}{\ensuremath{^{1}}}
\DeclareUnicodeCharacter{00B2}{\ensuremath{^{2}}}
\DeclareUnicodeCharacter{00B3}{\ensuremath{^{3}}}
\DeclareUnicodeCharacter{2074}{\ensuremath{^{4}}}
\DeclareUnicodeCharacter{2075}{\ensuremath{^{5}}}
\DeclareUnicodeCharacter{2076}{\ensuremath{^{6}}}
\DeclareUnicodeCharacter{2077}{\ensuremath{^{7}}}
\DeclareUnicodeCharacter{2078}{\ensuremath{^{8}}}
\DeclareUnicodeCharacter{2079}{\ensuremath{^{9}}}
\DeclareUnicodeCharacter{2071}{\ensuremath{^{i}}}
\DeclareUnicodeCharacter{2080}{\ensuremath{_{0}}}
\DeclareUnicodeCharacter{2081}{\ensuremath{_{1}}}
\DeclareUnicodeCharacter{2082}{\ensuremath{_{2}}}
\DeclareUnicodeCharacter{2083}{\ensuremath{_{3}}}
\DeclareUnicodeCharacter{2084}{\ensuremath{_{4}}}
\DeclareUnicodeCharacter{2085}{\ensuremath{_{5}}}
\DeclareUnicodeCharacter{2086}{\ensuremath{_{6}}}
\DeclareUnicodeCharacter{2087}{\ensuremath{_{7}}}
\DeclareUnicodeCharacter{2088}{\ensuremath{_{8}}}
\DeclareUnicodeCharacter{2089}{\ensuremath{_{9}}}
\DeclareUnicodeCharacter{208A}{\ensuremath{_{+}}}
\DeclareUnicodeCharacter{208B}{\ensuremath{_{-}}}
\DeclareUnicodeCharacter{2192}{\ensuremath{\to}}
\DeclareUnicodeCharacter{2502}{|}
\DeclareUnicodeCharacter{25BC}{\ensuremath{\blacktriangledown}}
\DeclareUnicodeCharacter{2500}{\ensuremath{-}}
\DeclareUnicodeCharacter{2514}{\ensuremath{\llcorner}}
\DeclareUnicodeCharacter{251C}{\ensuremath{\vdash}}
\DeclareUnicodeCharacter{2026}{\ldots}
\DeclareUnicodeCharacter{2208}{\ensuremath{\in}}
\DeclareUnicodeCharacter{2229}{\ensuremath{\cap}}
\DeclareUnicodeCharacter{222A}{\ensuremath{\cup}}
\DeclareUnicodeCharacter{2295}{\ensuremath{\oplus}}
\DeclareUnicodeCharacter{2297}{\ensuremath{\otimes}}
\DeclareUnicodeCharacter{2245}{\ensuremath{\cong}}
\DeclareUnicodeCharacter{2248}{\ensuremath{\approx}}
\DeclareUnicodeCharacter{2260}{\ensuremath{\neq}}
\DeclareUnicodeCharacter{2261}{\ensuremath{\equiv}}
\DeclareUnicodeCharacter{2264}{\ensuremath{\leq}}
\DeclareUnicodeCharacter{2265}{\ensuremath{\geq}}
\DeclareUnicodeCharacter{00D7}{\ensuremath{\times}}
\DeclareUnicodeCharacter{00B1}{\ensuremath{\pm}}
\DeclareUnicodeCharacter{2225}{\ensuremath{\|}}
% Greek letters used in body text (pandoc passes them through unchanged)
\DeclareUnicodeCharacter{03B1}{\ensuremath{\alpha}}
\DeclareUnicodeCharacter{03B2}{\ensuremath{\beta}}
\DeclareUnicodeCharacter{03B3}{\ensuremath{\gamma}}
\DeclareUnicodeCharacter{03B4}{\ensuremath{\delta}}
\DeclareUnicodeCharacter{03B5}{\ensuremath{\varepsilon}}
\DeclareUnicodeCharacter{03B6}{\ensuremath{\zeta}}
\DeclareUnicodeCharacter{03B7}{\ensuremath{\eta}}
\DeclareUnicodeCharacter{03B8}{\ensuremath{\theta}}
\DeclareUnicodeCharacter{03B9}{\ensuremath{\iota}}
\DeclareUnicodeCharacter{03BA}{\ensuremath{\kappa}}
\DeclareUnicodeCharacter{03BB}{\ensuremath{\lambda}}
\DeclareUnicodeCharacter{03BC}{\ensuremath{\mu}}
\DeclareUnicodeCharacter{03BD}{\ensuremath{\nu}}
\DeclareUnicodeCharacter{03BE}{\ensuremath{\xi}}
\DeclareUnicodeCharacter{03C0}{\ensuremath{\pi}}
\DeclareUnicodeCharacter{03C1}{\ensuremath{\rho}}
\DeclareUnicodeCharacter{03C3}{\ensuremath{\sigma}}
\DeclareUnicodeCharacter{03C4}{\ensuremath{\tau}}
\DeclareUnicodeCharacter{03C5}{\ensuremath{\upsilon}}
\DeclareUnicodeCharacter{03C6}{\ensuremath{\varphi}}
\DeclareUnicodeCharacter{03C7}{\ensuremath{\chi}}
\DeclareUnicodeCharacter{03C8}{\ensuremath{\psi}}
\DeclareUnicodeCharacter{03C9}{\ensuremath{\omega}}
\DeclareUnicodeCharacter{0394}{\ensuremath{\Delta}}
\DeclareUnicodeCharacter{0398}{\ensuremath{\Theta}}
\DeclareUnicodeCharacter{039B}{\ensuremath{\Lambda}}
\DeclareUnicodeCharacter{03A0}{\ensuremath{\Pi}}
\DeclareUnicodeCharacter{03A3}{\ensuremath{\Sigma}}
\DeclareUnicodeCharacter{03A6}{\ensuremath{\Phi}}
\DeclareUnicodeCharacter{03A8}{\ensuremath{\Psi}}
\DeclareUnicodeCharacter{03A9}{\ensuremath{\Omega}}
\DeclareUnicodeCharacter{2200}{\ensuremath{\forall}}
\DeclareUnicodeCharacter{2203}{\ensuremath{\exists}}
\DeclareUnicodeCharacter{2207}{\ensuremath{\nabla}}
\DeclareUnicodeCharacter{2202}{\ensuremath{\partial}}
\DeclareUnicodeCharacter{2218}{\ensuremath{\circ}}
\DeclareUnicodeCharacter{2032}{\ensuremath{\prime}}
\DeclareUnicodeCharacter{2113}{\ensuremath{\ell}}
\DeclareUnicodeCharacter{210F}{\ensuremath{\hbar}}
\DeclareUnicodeCharacter{27E8}{\ensuremath{\langle}}
\DeclareUnicodeCharacter{27E9}{\ensuremath{\rangle}}
\DeclareUnicodeCharacter{2308}{\ensuremath{\lceil}}
\DeclareUnicodeCharacter{2309}{\ensuremath{\rceil}}
\DeclareUnicodeCharacter{230A}{\ensuremath{\lfloor}}
\DeclareUnicodeCharacter{230B}{\ensuremath{\rfloor}}
\DeclareUnicodeCharacter{222B}{\ensuremath{\int}}
\DeclareUnicodeCharacter{2211}{\ensuremath{\sum}}
\DeclareUnicodeCharacter{220F}{\ensuremath{\prod}}
\DeclareUnicodeCharacter{221E}{\ensuremath{\infty}}
\DeclareUnicodeCharacter{00B5}{\ensuremath{\mu}}
\DeclareUnicodeCharacter{2236}{\ensuremath{\colon}}
\DeclareUnicodeCharacter{2192}{\ensuremath{\to}}
\DeclareUnicodeCharacter{21A6}{\ensuremath{\mapsto}}
\DeclareUnicodeCharacter{21D2}{\ensuremath{\Rightarrow}}
\DeclareUnicodeCharacter{21D4}{\ensuremath{\Leftrightarrow}}
\DeclareUnicodeCharacter{2282}{\ensuremath{\subset}}
\DeclareUnicodeCharacter{2286}{\ensuremath{\subseteq}}
\DeclareUnicodeCharacter{2287}{\ensuremath{\supseteq}}
\DeclareUnicodeCharacter{222B}{\ensuremath{\int}}
\DeclareUnicodeCharacter{2205}{\ensuremath{\emptyset}}
\DeclareUnicodeCharacter{2299}{\ensuremath{\odot}}
\DeclareUnicodeCharacter{2220}{\ensuremath{\angle}}
\DeclareUnicodeCharacter{2135}{\ensuremath{\aleph}}
\DeclareUnicodeCharacter{207B}{\ensuremath{^{-}}}
\DeclareUnicodeCharacter{207A}{\ensuremath{^{+}}}
\DeclareUnicodeCharacter{220E}{\ensuremath{\blacksquare}}
\DeclareUnicodeCharacter{2061}{}
\DeclareUnicodeCharacter{2062}{}
\DeclareUnicodeCharacter{2063}{}
\DeclareUnicodeCharacter{2009}{\,}
\DeclareUnicodeCharacter{200A}{\,}
\DeclareUnicodeCharacter{200B}{}
\DeclareUnicodeCharacter{2212}{\ensuremath{-}}
\DeclareUnicodeCharacter{2194}{\ensuremath{\leftrightarrow}}
\DeclareUnicodeCharacter{2099}{\ensuremath{_{n}}}
\DeclareUnicodeCharacter{2098}{\ensuremath{_{m}}}
\DeclareUnicodeCharacter{2096}{\ensuremath{_{k}}}
\DeclareUnicodeCharacter{2095}{\ensuremath{_{h}}}
\DeclareUnicodeCharacter{2090}{\ensuremath{_{a}}}
\DeclareUnicodeCharacter{2091}{\ensuremath{_{e}}}
\DeclareUnicodeCharacter{2093}{\ensuremath{_{x}}}
\DeclareUnicodeCharacter{1D62}{\ensuremath{_{i}}}
\DeclareUnicodeCharacter{2C7C}{\ensuremath{_{j}}}
\DeclareUnicodeCharacter{2C7B}{\ensuremath{^{e}}}
\DeclareUnicodeCharacter{1D52}{\ensuremath{^{o}}}
\DeclareUnicodeCharacter{1D5}{\ensuremath{^{u}}}
\DeclareUnicodeCharacter{02E1}{\ensuremath{^{l}}}
\DeclareUnicodeCharacter{2295}{\ensuremath{\oplus}}
\DeclareUnicodeCharacter{27FA}{\ensuremath{\Longleftrightarrow}}
\DeclareUnicodeCharacter{27F9}{\ensuremath{\Longrightarrow}}
\DeclareUnicodeCharacter{2299}{\ensuremath{\odot}}
% FoP-specific symbols (blackboard division algebras, set/relation operators)
\DeclareUnicodeCharacter{00F7}{\ensuremath{\div}}
\DeclareUnicodeCharacter{2102}{\ensuremath{\mathbb{C}}}
\DeclareUnicodeCharacter{210D}{\ensuremath{\mathbb{H}}}
\DeclareUnicodeCharacter{211D}{\ensuremath{\mathbb{R}}}
\DeclareUnicodeCharacter{1D546}{\ensuremath{\mathbb{O}}}
\DeclareUnicodeCharacter{2227}{\ensuremath{\wedge}}
\DeclareUnicodeCharacter{2272}{\ensuremath{\lesssim}}
\DeclareUnicodeCharacter{2283}{\ensuremath{\supset}}
\DeclareUnicodeCharacter{228A}{\ensuremath{\subsetneq}}
\DeclareUnicodeCharacter{2016}{\ensuremath{\|}}
\DeclareUnicodeCharacter{2190}{\ensuremath{\leftarrow}}
\DeclareUnicodeCharacter{2228}{\ensuremath{\vee}}
\DeclareUnicodeCharacter{22C6}{\ensuremath{\star}}
\DeclareUnicodeCharacter{1D4D8}{\ensuremath{\mathcal{I}}}
\DeclareUnicodeCharacter{1D524}{\ensuremath{\mathfrak{g}}}
\DeclareUnicodeCharacter{23F3}{\textsf{[open]}}
% Provide \tightlist used by pandoc-generated lists (no-op)
\providecommand{\tightlist}{\setlength{\itemsep}{0pt}\setlength{\parskip}{0pt}}
% pandoc --listings emits \passthrough around inline \lstinline spans
\providecommand{\passthrough}[1]{#1}
% pandoc longtable column specs: provide \real as passthrough and a
% `none` counter so that calc's \real{x} expansion path works.
\newcounter{none}
\providecommand{\real}[1]{#1}
% pandoc >=3.x wraps figures in \pandocbounded (official definition: scale to
% fit text width/height, never upscale). Requires graphicx (loaded above).
\makeatletter
\newsavebox\pandoc@box
\providecommand*\pandocbounded[1]{%
  \sbox\pandoc@box{#1}%
  \ifdim\wd\pandoc@box>\linewidth
    \resizebox{\linewidth}{!}{\usebox\pandoc@box}%
  \else
    \usebox\pandoc@box
  \fi}
\makeatother

% LISTINGS
\lstset{
    basicstyle=\small\ttfamily,
    breaklines=true, frame=single,
    keepspaces=true, showstringspaces=false,
    aboveskip=0.8em, belowskip=0.8em
}

% HEADER/FOOTER
\setlength{\headheight}{14pt}
\pagestyle{fancy}
\fancyhf{}
\fancyhead[R]{\thepage}
\renewcommand{\headrulewidth}{0.2pt}

% HYPERREF
\hypersetup{
    colorlinks=true,
    linkcolor=blue, citecolor=blue, urlcolor=blue,
    pdfauthor={Brieuc de La Fourniere}
}

% SPACING
\setstretch{1.2}
\setlength{\parskip}{0.4em}
\setlength{\parindent}{0pt}

% TITLE FORMATTING
\usepackage{titling}
\pretitle{\LARGE\bfseries}
\posttitle{\vspace{-0.4em}}
\preauthor{}
\postauthor{}
\predate{}
\postdate{}
\setlength{\droptitle}{-2.0em}

% CUSTOM COMMANDS
\newcommand{\E}{\mathrm{E}}
\newcommand{\Gtwo}{\mathrm{G}_2}
\newcommand{\Kseven}{K_7}
\newcommand{\AdS}{\mathrm{AdS}}
\newcommand{\SU}{\mathrm{SU}}
\newcommand{\SO}{\mathrm{SO}}
\newcommand{\U}{\mathrm{U}}
\newcommand{\Sp}{\mathrm{Sp}}
\newcommand{\PSL}{\mathrm{PSL}}
\newcommand{\SM}{\mathrm{SM}}
\newcommand{\Tr}{\mathrm{Tr}}
\newcommand{\rank}{\mathrm{rank}}
\newcommand{\Vol}{\mathrm{Vol}}
\newcommand{\Hol}{\mathrm{Hol}}
\newcommand{\Ric}{\mathrm{Ric}}
\newcommand{\Rm}{\mathrm{Rm}}
\newcommand{\Scal}{\mathrm{Scal}}
\DeclareMathOperator{\diag}{diag}
\DeclareMathOperator{\cond}{cond}
\newcommand{\CP}{\mathrm{CP}}
\newcommand{\CKM}{\mathrm{CKM}}
\newcommand{\PMNS}{\mathrm{PMNS}}
\newcommand{\EW}{\mathrm{EW}}
\newcommand{\Pl}{\mathrm{Pl}}
\newcommand{\DE}{\mathrm{DE}}
\newcommand{\DM}{\mathrm{DM}}
\newcommand{\KK}{\mathrm{KK}}
\newcommand{\NK}{\mathrm{NK}}
\newcommand{\TCS}{\mathrm{TCS}}
\newcommand{\PINN}{\mathrm{PINN}}
\newcommand{\btwo}{b_2}
\newcommand{\bthree}{b_3}
\newcommand{\typeI}{\textsc{I}}
\newcommand{\typeII}{\textsc{II}}
\newcommand{\typeIII}{\textsc{III}}
\newcommand{\typeIV}{\textsc{IV}}
\newcommand{\proven}{\textsc{Proven}}

% PDF bookmark safety
\pdfstringdefDisableCommands{%
  \def\textsubscript#1{#1}%
  \def\textsuperscript#1{#1}%
  \def\CP{CP}%
  \def\Gtwo{G2}%
  \def\E{E}%
  \def\SM{SM}%
}

% THEOREM ENVIRONMENTS
\newtheorem{theorem}{Theorem}[section]
\newtheorem{proposition}[theorem]{Proposition}
\newtheorem{lemma}[theorem]{Lemma}
\newtheorem{corollary}[theorem]{Corollary}
\theoremstyle{definition}
\newtheorem{definition}[theorem]{Definition}
\theoremstyle{remark}
\newtheorem{remark}[theorem]{Remark}
 after \documentclass

% ENCODING & FONTS
\usepackage[utf8]{inputenc}
\usepackage[T1]{fontenc}
\usepackage{lmodern}

% PAGE LAYOUT (Golden Ratio)
\usepackage[margin=1.618cm, top=2.618cm, bottom=2.618cm]{geometry}

% ESSENTIAL PACKAGES
\usepackage{float}
\usepackage{caption}
\usepackage{subcaption}
\usepackage{setspace}
\usepackage{fancyhdr}
\usepackage{xcolor}
\usepackage{hyperref}
\usepackage{csquotes}
\usepackage{amsmath}
\usepackage{amssymb}
\usepackage{amsthm}
\usepackage{mathtools}
\usepackage{booktabs}
\usepackage{longtable}
\usepackage{array}
\usepackage{multirow}
\usepackage{tikz}
\usepackage{graphicx}
\usepackage{listings}
\usepackage{enumitem}
\usepackage{etoolbox}
% Render generated tables a notch smaller so wide pipe-tables breathe.
\AtBeginEnvironment{longtable}{\footnotesize}
\AtBeginEnvironment{tabular}{\small}
% Unicode fallbacks (pandoc artifacts that survive markdown→LaTeX)
\DeclareUnicodeCharacter{00B0}{\ensuremath{^\circ}}
\DeclareUnicodeCharacter{221A}{\ensuremath{\sqrt{\,}}}
\DeclareUnicodeCharacter{2713}{\ensuremath{\checkmark}}
\DeclareUnicodeCharacter{2717}{\ensuremath{\times}}
\DeclareUnicodeCharacter{2070}{\ensuremath{^{0}}}
\DeclareUnicodeCharacter{00B9}{\ensuremath{^{1}}}
\DeclareUnicodeCharacter{00B2}{\ensuremath{^{2}}}
\DeclareUnicodeCharacter{00B3}{\ensuremath{^{3}}}
\DeclareUnicodeCharacter{2074}{\ensuremath{^{4}}}
\DeclareUnicodeCharacter{2075}{\ensuremath{^{5}}}
\DeclareUnicodeCharacter{2076}{\ensuremath{^{6}}}
\DeclareUnicodeCharacter{2077}{\ensuremath{^{7}}}
\DeclareUnicodeCharacter{2078}{\ensuremath{^{8}}}
\DeclareUnicodeCharacter{2079}{\ensuremath{^{9}}}
\DeclareUnicodeCharacter{2071}{\ensuremath{^{i}}}
\DeclareUnicodeCharacter{2080}{\ensuremath{_{0}}}
\DeclareUnicodeCharacter{2081}{\ensuremath{_{1}}}
\DeclareUnicodeCharacter{2082}{\ensuremath{_{2}}}
\DeclareUnicodeCharacter{2083}{\ensuremath{_{3}}}
\DeclareUnicodeCharacter{2084}{\ensuremath{_{4}}}
\DeclareUnicodeCharacter{2085}{\ensuremath{_{5}}}
\DeclareUnicodeCharacter{2086}{\ensuremath{_{6}}}
\DeclareUnicodeCharacter{2087}{\ensuremath{_{7}}}
\DeclareUnicodeCharacter{2088}{\ensuremath{_{8}}}
\DeclareUnicodeCharacter{2089}{\ensuremath{_{9}}}
\DeclareUnicodeCharacter{208A}{\ensuremath{_{+}}}
\DeclareUnicodeCharacter{208B}{\ensuremath{_{-}}}
\DeclareUnicodeCharacter{2192}{\ensuremath{\to}}
\DeclareUnicodeCharacter{2502}{|}
\DeclareUnicodeCharacter{25BC}{\ensuremath{\blacktriangledown}}
\DeclareUnicodeCharacter{2500}{\ensuremath{-}}
\DeclareUnicodeCharacter{2514}{\ensuremath{\llcorner}}
\DeclareUnicodeCharacter{251C}{\ensuremath{\vdash}}
\DeclareUnicodeCharacter{2026}{\ldots}
\DeclareUnicodeCharacter{2208}{\ensuremath{\in}}
\DeclareUnicodeCharacter{2229}{\ensuremath{\cap}}
\DeclareUnicodeCharacter{222A}{\ensuremath{\cup}}
\DeclareUnicodeCharacter{2295}{\ensuremath{\oplus}}
\DeclareUnicodeCharacter{2297}{\ensuremath{\otimes}}
\DeclareUnicodeCharacter{2245}{\ensuremath{\cong}}
\DeclareUnicodeCharacter{2248}{\ensuremath{\approx}}
\DeclareUnicodeCharacter{2260}{\ensuremath{\neq}}
\DeclareUnicodeCharacter{2261}{\ensuremath{\equiv}}
\DeclareUnicodeCharacter{2264}{\ensuremath{\leq}}
\DeclareUnicodeCharacter{2265}{\ensuremath{\geq}}
\DeclareUnicodeCharacter{00D7}{\ensuremath{\times}}
\DeclareUnicodeCharacter{00B1}{\ensuremath{\pm}}
\DeclareUnicodeCharacter{2225}{\ensuremath{\|}}
% Greek letters used in body text (pandoc passes them through unchanged)
\DeclareUnicodeCharacter{03B1}{\ensuremath{\alpha}}
\DeclareUnicodeCharacter{03B2}{\ensuremath{\beta}}
\DeclareUnicodeCharacter{03B3}{\ensuremath{\gamma}}
\DeclareUnicodeCharacter{03B4}{\ensuremath{\delta}}
\DeclareUnicodeCharacter{03B5}{\ensuremath{\varepsilon}}
\DeclareUnicodeCharacter{03B6}{\ensuremath{\zeta}}
\DeclareUnicodeCharacter{03B7}{\ensuremath{\eta}}
\DeclareUnicodeCharacter{03B8}{\ensuremath{\theta}}
\DeclareUnicodeCharacter{03B9}{\ensuremath{\iota}}
\DeclareUnicodeCharacter{03BA}{\ensuremath{\kappa}}
\DeclareUnicodeCharacter{03BB}{\ensuremath{\lambda}}
\DeclareUnicodeCharacter{03BC}{\ensuremath{\mu}}
\DeclareUnicodeCharacter{03BD}{\ensuremath{\nu}}
\DeclareUnicodeCharacter{03BE}{\ensuremath{\xi}}
\DeclareUnicodeCharacter{03C0}{\ensuremath{\pi}}
\DeclareUnicodeCharacter{03C1}{\ensuremath{\rho}}
\DeclareUnicodeCharacter{03C3}{\ensuremath{\sigma}}
\DeclareUnicodeCharacter{03C4}{\ensuremath{\tau}}
\DeclareUnicodeCharacter{03C5}{\ensuremath{\upsilon}}
\DeclareUnicodeCharacter{03C6}{\ensuremath{\varphi}}
\DeclareUnicodeCharacter{03C7}{\ensuremath{\chi}}
\DeclareUnicodeCharacter{03C8}{\ensuremath{\psi}}
\DeclareUnicodeCharacter{03C9}{\ensuremath{\omega}}
\DeclareUnicodeCharacter{0394}{\ensuremath{\Delta}}
\DeclareUnicodeCharacter{0398}{\ensuremath{\Theta}}
\DeclareUnicodeCharacter{039B}{\ensuremath{\Lambda}}
\DeclareUnicodeCharacter{03A0}{\ensuremath{\Pi}}
\DeclareUnicodeCharacter{03A3}{\ensuremath{\Sigma}}
\DeclareUnicodeCharacter{03A6}{\ensuremath{\Phi}}
\DeclareUnicodeCharacter{03A8}{\ensuremath{\Psi}}
\DeclareUnicodeCharacter{03A9}{\ensuremath{\Omega}}
\DeclareUnicodeCharacter{2200}{\ensuremath{\forall}}
\DeclareUnicodeCharacter{2203}{\ensuremath{\exists}}
\DeclareUnicodeCharacter{2207}{\ensuremath{\nabla}}
\DeclareUnicodeCharacter{2202}{\ensuremath{\partial}}
\DeclareUnicodeCharacter{2218}{\ensuremath{\circ}}
\DeclareUnicodeCharacter{2032}{\ensuremath{\prime}}
\DeclareUnicodeCharacter{2113}{\ensuremath{\ell}}
\DeclareUnicodeCharacter{210F}{\ensuremath{\hbar}}
\DeclareUnicodeCharacter{27E8}{\ensuremath{\langle}}
\DeclareUnicodeCharacter{27E9}{\ensuremath{\rangle}}
\DeclareUnicodeCharacter{2308}{\ensuremath{\lceil}}
\DeclareUnicodeCharacter{2309}{\ensuremath{\rceil}}
\DeclareUnicodeCharacter{230A}{\ensuremath{\lfloor}}
\DeclareUnicodeCharacter{230B}{\ensuremath{\rfloor}}
\DeclareUnicodeCharacter{222B}{\ensuremath{\int}}
\DeclareUnicodeCharacter{2211}{\ensuremath{\sum}}
\DeclareUnicodeCharacter{220F}{\ensuremath{\prod}}
\DeclareUnicodeCharacter{221E}{\ensuremath{\infty}}
\DeclareUnicodeCharacter{00B5}{\ensuremath{\mu}}
\DeclareUnicodeCharacter{2236}{\ensuremath{\colon}}
\DeclareUnicodeCharacter{2192}{\ensuremath{\to}}
\DeclareUnicodeCharacter{21A6}{\ensuremath{\mapsto}}
\DeclareUnicodeCharacter{21D2}{\ensuremath{\Rightarrow}}
\DeclareUnicodeCharacter{21D4}{\ensuremath{\Leftrightarrow}}
\DeclareUnicodeCharacter{2282}{\ensuremath{\subset}}
\DeclareUnicodeCharacter{2286}{\ensuremath{\subseteq}}
\DeclareUnicodeCharacter{2287}{\ensuremath{\supseteq}}
\DeclareUnicodeCharacter{222B}{\ensuremath{\int}}
\DeclareUnicodeCharacter{2205}{\ensuremath{\emptyset}}
\DeclareUnicodeCharacter{2299}{\ensuremath{\odot}}
\DeclareUnicodeCharacter{2220}{\ensuremath{\angle}}
\DeclareUnicodeCharacter{2135}{\ensuremath{\aleph}}
\DeclareUnicodeCharacter{207B}{\ensuremath{^{-}}}
\DeclareUnicodeCharacter{207A}{\ensuremath{^{+}}}
\DeclareUnicodeCharacter{220E}{\ensuremath{\blacksquare}}
\DeclareUnicodeCharacter{2061}{}
\DeclareUnicodeCharacter{2062}{}
\DeclareUnicodeCharacter{2063}{}
\DeclareUnicodeCharacter{2009}{\,}
\DeclareUnicodeCharacter{200A}{\,}
\DeclareUnicodeCharacter{200B}{}
\DeclareUnicodeCharacter{2212}{\ensuremath{-}}
\DeclareUnicodeCharacter{2194}{\ensuremath{\leftrightarrow}}
\DeclareUnicodeCharacter{2099}{\ensuremath{_{n}}}
\DeclareUnicodeCharacter{2098}{\ensuremath{_{m}}}
\DeclareUnicodeCharacter{2096}{\ensuremath{_{k}}}
\DeclareUnicodeCharacter{2095}{\ensuremath{_{h}}}
\DeclareUnicodeCharacter{2090}{\ensuremath{_{a}}}
\DeclareUnicodeCharacter{2091}{\ensuremath{_{e}}}
\DeclareUnicodeCharacter{2093}{\ensuremath{_{x}}}
\DeclareUnicodeCharacter{1D62}{\ensuremath{_{i}}}
\DeclareUnicodeCharacter{2C7C}{\ensuremath{_{j}}}
\DeclareUnicodeCharacter{2C7B}{\ensuremath{^{e}}}
\DeclareUnicodeCharacter{1D52}{\ensuremath{^{o}}}
\DeclareUnicodeCharacter{1D5}{\ensuremath{^{u}}}
\DeclareUnicodeCharacter{02E1}{\ensuremath{^{l}}}
\DeclareUnicodeCharacter{2295}{\ensuremath{\oplus}}
\DeclareUnicodeCharacter{27FA}{\ensuremath{\Longleftrightarrow}}
\DeclareUnicodeCharacter{27F9}{\ensuremath{\Longrightarrow}}
\DeclareUnicodeCharacter{2299}{\ensuremath{\odot}}
% FoP-specific symbols (blackboard division algebras, set/relation operators)
\DeclareUnicodeCharacter{00F7}{\ensuremath{\div}}
\DeclareUnicodeCharacter{2102}{\ensuremath{\mathbb{C}}}
\DeclareUnicodeCharacter{210D}{\ensuremath{\mathbb{H}}}
\DeclareUnicodeCharacter{211D}{\ensuremath{\mathbb{R}}}
\DeclareUnicodeCharacter{1D546}{\ensuremath{\mathbb{O}}}
\DeclareUnicodeCharacter{2227}{\ensuremath{\wedge}}
\DeclareUnicodeCharacter{2272}{\ensuremath{\lesssim}}
\DeclareUnicodeCharacter{2283}{\ensuremath{\supset}}
\DeclareUnicodeCharacter{228A}{\ensuremath{\subsetneq}}
\DeclareUnicodeCharacter{2016}{\ensuremath{\|}}
\DeclareUnicodeCharacter{2190}{\ensuremath{\leftarrow}}
\DeclareUnicodeCharacter{2228}{\ensuremath{\vee}}
\DeclareUnicodeCharacter{22C6}{\ensuremath{\star}}
\DeclareUnicodeCharacter{1D4D8}{\ensuremath{\mathcal{I}}}
\DeclareUnicodeCharacter{1D524}{\ensuremath{\mathfrak{g}}}
\DeclareUnicodeCharacter{23F3}{\textsf{[open]}}
% Provide \tightlist used by pandoc-generated lists (no-op)
\providecommand{\tightlist}{\setlength{\itemsep}{0pt}\setlength{\parskip}{0pt}}
% pandoc --listings emits \passthrough around inline \lstinline spans
\providecommand{\passthrough}[1]{#1}
% pandoc longtable column specs: provide \real as passthrough and a
% `none` counter so that calc's \real{x} expansion path works.
\newcounter{none}
\providecommand{\real}[1]{#1}
% pandoc >=3.x wraps figures in \pandocbounded (official definition: scale to
% fit text width/height, never upscale). Requires graphicx (loaded above).
\makeatletter
\newsavebox\pandoc@box
\providecommand*\pandocbounded[1]{%
  \sbox\pandoc@box{#1}%
  \ifdim\wd\pandoc@box>\linewidth
    \resizebox{\linewidth}{!}{\usebox\pandoc@box}%
  \else
    \usebox\pandoc@box
  \fi}
\makeatother

% LISTINGS
\lstset{
    basicstyle=\small\ttfamily,
    breaklines=true, frame=single,
    keepspaces=true, showstringspaces=false,
    aboveskip=0.8em, belowskip=0.8em
}

% HEADER/FOOTER
\setlength{\headheight}{14pt}
\pagestyle{fancy}
\fancyhf{}
\fancyhead[R]{\thepage}
\renewcommand{\headrulewidth}{0.2pt}

% HYPERREF
\hypersetup{
    colorlinks=true,
    linkcolor=blue, citecolor=blue, urlcolor=blue,
    pdfauthor={Brieuc de La Fourniere}
}

% SPACING
\setstretch{1.2}
\setlength{\parskip}{0.4em}
\setlength{\parindent}{0pt}

% TITLE FORMATTING
\usepackage{titling}
\pretitle{\LARGE\bfseries}
\posttitle{\vspace{-0.4em}}
\preauthor{}
\postauthor{}
\predate{}
\postdate{}
\setlength{\droptitle}{-2.0em}

% CUSTOM COMMANDS
\newcommand{\E}{\mathrm{E}}
\newcommand{\Gtwo}{\mathrm{G}_2}
\newcommand{\Kseven}{K_7}
\newcommand{\AdS}{\mathrm{AdS}}
\newcommand{\SU}{\mathrm{SU}}
\newcommand{\SO}{\mathrm{SO}}
\newcommand{\U}{\mathrm{U}}
\newcommand{\Sp}{\mathrm{Sp}}
\newcommand{\PSL}{\mathrm{PSL}}
\newcommand{\SM}{\mathrm{SM}}
\newcommand{\Tr}{\mathrm{Tr}}
\newcommand{\rank}{\mathrm{rank}}
\newcommand{\Vol}{\mathrm{Vol}}
\newcommand{\Hol}{\mathrm{Hol}}
\newcommand{\Ric}{\mathrm{Ric}}
\newcommand{\Rm}{\mathrm{Rm}}
\newcommand{\Scal}{\mathrm{Scal}}
\DeclareMathOperator{\diag}{diag}
\DeclareMathOperator{\cond}{cond}
\newcommand{\CP}{\mathrm{CP}}
\newcommand{\CKM}{\mathrm{CKM}}
\newcommand{\PMNS}{\mathrm{PMNS}}
\newcommand{\EW}{\mathrm{EW}}
\newcommand{\Pl}{\mathrm{Pl}}
\newcommand{\DE}{\mathrm{DE}}
\newcommand{\DM}{\mathrm{DM}}
\newcommand{\KK}{\mathrm{KK}}
\newcommand{\NK}{\mathrm{NK}}
\newcommand{\TCS}{\mathrm{TCS}}
\newcommand{\PINN}{\mathrm{PINN}}
\newcommand{\btwo}{b_2}
\newcommand{\bthree}{b_3}
\newcommand{\typeI}{\textsc{I}}
\newcommand{\typeII}{\textsc{II}}
\newcommand{\typeIII}{\textsc{III}}
\newcommand{\typeIV}{\textsc{IV}}
\newcommand{\proven}{\textsc{Proven}}

% PDF bookmark safety
\pdfstringdefDisableCommands{%
  \def\textsubscript#1{#1}%
  \def\textsuperscript#1{#1}%
  \def\CP{CP}%
  \def\Gtwo{G2}%
  \def\E{E}%
  \def\SM{SM}%
}

% THEOREM ENVIRONMENTS
\newtheorem{theorem}{Theorem}[section]
\newtheorem{proposition}[theorem]{Proposition}
\newtheorem{lemma}[theorem]{Lemma}
\newtheorem{corollary}[theorem]{Corollary}
\theoremstyle{definition}
\newtheorem{definition}[theorem]{Definition}
\theoremstyle{remark}
\newtheorem{remark}[theorem]{Remark}
